\chapter{Introduction} \label{cha:introduction}
\pagestyle{plain}

Large language models (LLMs) process text as sequences of discrete tokens rather than whole words \citep{mielke2021}, which can split meaningful units into smaller fragments. For example, the GPT-2 tokenizer represents ``Hemingway'' as \texttt{['␣H', 'eming', 'way']}\footnote{\texttt{␣} refers to whitespace, since many words in English tokenize differently depending on if the word has a space in front of it or not.}. The choice in tokenization scheme has been shown to have impact on downstream performance, especially in lower-resource languages and languages with non-Latin scripts \citep{ahia2023, petrov2023}. Despite this, models remain robust to input corruption \citep{cao2023}, raising the question of how models recover coherent semantic units. One hypothesis, often called \textit{detokenization} \citep{lad2024}, suggests that early layers partially reverse tokenization by combining tokens into larger segments. Interpretability studies \citep{elhage2022a} provide supporting evidence, identifying components that respond to token boundary artifacts in early layers.

\citet{kaplan2025} describe detokenization as involving two main components: early-layer attention mechanisms that group tokens, and feedforward networks in early–mid layers that refine these into more coherent representations. They further propose that LLMs develop an \textit{inner lexicon} that extends beyond the tokenizer, with detokenization effectively being an attempt at mapping the input embeddings/positional encodings to this internal representation for downstream processing.

This thesis investigates whether detokenization depends on a specific set of attention heads and how it varies across languages (Chinese, Turkish, and English), tokenization schemes, and model sizes. Two dimensions are examined: differences between otherwise comparable models using BPE-, SentencePiece-, and byte-level tokenization, and differences in scale within a single model family. Using activation patching \citep{vig2019, ghandeharioun2024}, the study evaluates performance on a word retrieval task in which a model reconstructs a full word from the representation of its final token. Retrieval performance is probed across layers under original, targeted ablation (1/3/5\%), and random ablation conditions. This setup enables analysis of whether detokenization relies on specific attention heads or whether later components compensate for earlier disruptions. The thesis aims to clarify how tokenizer design and multilingual input affect the formation of intermediate word representations in LLMs.

Section \ref{cha:background} provides background, Section \ref{cha:method} outlines the approach, Section \ref{cha:results} presents the results, and Section \ref{cha:discussion} discusses them.

\section*{Research Questions}
\begin{description}
\item[\textbf{RQ1}]\label{rq1} Is there a specific set of attention heads specialized to form these intermediate representations?
\item[\textbf{RQ2}]\label{rq2} Does the choice of tokenizer play a role in how well the intermediate representation are formed and how robust the model is to attention head ablation?
\item[\textbf{RQ3}]\label{rq3} Does the parameter size play a role in how well the intermediate representation is formed and how robust the model is to attention head ablation?
\end{description}

\section*{Collaboration}

The thesis is done in collaboration with the NLP group at RISE.
